\section{Review Method}\label{sec:method}

The review rests on the Networked XR Atlas, whose full protocol, code and intermediate decisions are
public. Here we summarise the steps that bear on how the evidence should be read. \Cref{fig:prisma} shows
the flow of records.

\paragraph{Research questions.} (RQ1) How large is the networked-XR literature, and how has its
composition changed? (RQ2) What does it establish about each function of a networked XR system, and
with what evidence? (RQ3) How realistically are its claims tested? (RQ4) Which gaps recur across
functions, and which research would close them first?

\paragraph{Search.} We queried OpenAlex, which indexes journals, proceedings, theses, books and
preprints. A record qualified if its title or abstract contained one of \nxrQueryXrTerms{} XR phrases
\emph{and} one of \nxrQueryNetTerms{} networking phrases. The XR side covers virtual, augmented, mixed
and extended reality, the metaverse, 360-degree and volumetric video, holographic communication,
telepresence and the tactile internet. The networking side runs from streaming, edge computing and
remote rendering to 5G, 6G, Wi-Fi, QoE and caching. The XR side deliberately includes media terms. A
search for ``virtual reality'' alone would miss much of the field, because many papers describe their
content only as haptic traffic, omnidirectional video or point clouds. The window runs from October 2016
to October 2026. After removing \nxrDuplicates{} duplicates by identifier, DOI and title, \nxrAfterDedup{}
records remained.

\paragraph{Screening.} Records were sorted into tiers by how strongly XR and networking terms
co-occur. In a hand-labelled pilot of \nxrPilotN{} records, the weakest tier held no relevant paper, so
its \nxrTierCRule{} records were excluded by rule. The other \nxrJevScreened{} records were screened by a
calibrated classification model against a written scope. A paper was included when transmitting,
offloading, rendering remotely, caching, synchronising or securing XR content \emph{over a network} was
its core or a substantial subject. On the decidable pilot records, the inclusion rule reached precision
\nxrPilotPrecision{} and recall \nxrPilotRecall{}. A second de-duplication pass merged preprints with
their published versions, and retracted and off-topic records were removed (\nxrRemoved{} in all). The
final corpus holds \nxrIncluded{} papers: \nxrShareConfPapers{} conference papers, \nxrShareArticles{}
journal articles, and the rest preprints, theses and chapters.

\begin{figure}[t]\centering
\begin{tikzpicture}[font=\sffamily\footnotesize,
  box/.style={draw=black!45,rounded corners=2pt,align=center,minimum height=8.5mm,inner sep=3pt,fill=black!2},
  side/.style={draw=gapcolor!60,rounded corners=2pt,align=left,inner sep=3pt,fill=gapcolor!4,font=\sffamily\scriptsize},
  arr/.style={-{Stealth[length=2mm]},black!60,line width=0.6pt}]
  \node[box,text width=2.6cm] (h) {OpenAlex harvest\\\textbf{\nxrHarvested{}}};
  \node[box,text width=2.6cm,right=0.55cm of h] (d) {After de-duplication\\\textbf{\nxrAfterDedup{}}};
  \node[box,text width=2.6cm,right=0.55cm of d] (s) {Screened in\\\textbf{\nxrScreenedIncluded{}}};
  \node[box,text width=2.6cm,right=0.55cm of s,fill=encode!8,draw=encode!70] (i) {Included corpus\\\textbf{\nxrIncluded{}}};
  \draw[arr] (h)--(d); \draw[arr] (d)--(s); \draw[arr] (s)--(i);
  \node[side,below=0.35cm of d,text width=2.6cm] (x1) {Weak tier excluded by rule: \nxrTierCRule{}\\Model-excluded: \nxrExcluded{}\\Borderline: \nxrBorderline{}};
  \node[side,below=0.35cm of s,text width=2.6cm] (x2) {Duplicate versions, retracted, off-topic: \nxrRemoved{}};
  \draw[arr] (d.south)--(x1.north); \draw[arr] (s.south)--(x2.north);
  \node[box,below=0.35cm of i,text width=2.6cm] (l) {Labelled on 5 dimensions\\(all \nxrIncluded{})};
  \node[box,below=0.25cm of l,text width=2.6cm] (r) {Deep-read: \textbf{\nxrDeepReads{}} in \nxrDeepCells{} cells\\\nxrGaps{} gaps verified};
  \draw[arr] (i)--(l); \draw[arr] (l)--(r);
\end{tikzpicture}
\caption{Flow of records from harvest to analysis. Excluded-record counts are not additive with the
included count, because borderline records are kept apart and screening ran before the second
de-duplication pass.}\label{fig:prisma}
\end{figure}

\paragraph{Labelling.} Every included paper was labelled on five closed dimensions, each with written
definitions: the \emph{networking problem} it mainly addresses (\nxrNCatProb{} options), the \emph{network
technology} (\nxrNCatTech{}), the \emph{content type} transmitted (\nxrNCatCont{}), the \emph{method}
(\nxrNCatMeth{}) and the \emph{evaluation} (\nxrNCatEval{}). \nxrLabJev{} papers were labelled by the
model and the remaining \nxrLabDirect{} by direct reading of title and abstract. On \nxrRelabCal{} papers
labelled both ways, the primary labels agree for \nxrRelabAgrMin{}--\nxrRelabAgrMax{} of papers. The
top-two label sets overlap for \nxrRelabTwoMin{}--\nxrRelabTwoMax{}, so most disagreements concern papers
that sit between two categories. The labels support counts and trends over hundreds of papers. They do
not support claims about any single paper, and every paper discussed individually below was read.

\paragraph{Evaluation realism as evidence strength.} Reviews of interventions grade evidence by study
design. In networking the main threat to a claim is an unrealistic test environment, so we grade evidence
by \emph{evaluation realism}, from most to least realistic: a real network or testbed with human users; a
real network or testbed without users; emulation driven by real traces; simulation; analysis alone; and no
evaluation. Abstracts often blur ``prototype'' and ``emulation'', so this is the least reliable dimension
(top-1 agreement \nxrAccEval{} against \nxrAccCont{} for content). Differences of a few percentage points
in realism should not be over-read.

\paragraph{Deep reading and gap verification.} The atlas divides the corpus into \nxrDeepCells{} cells,
one per networking problem and one per specific content type. In each cell it selected the most
influential papers, by citations, citations per year and in-corpus PageRank, together with the strongest
recent papers and the main surveys. In all, \nxrDeepReads{} papers were read and summarised in a fixed
schema, \nxrDeepFulltext{} of them in full text. A candidate gap was kept only if it passed three tests.
It had to be meaningful rather than definitionally empty. It had to survive a regular-expression search
of every included abstract, which corrects for the single-label taxonomy. And the papers that had been
read had to leave it open. \nxrGaps{} gaps passed, \nxrGapsHigh{} of them with high confidence. The
open problems that close Sections~\ref{sec:encode}--\ref{sec:assure} are drawn from them.

\paragraph{Threats to validity.} OpenAlex lacks an abstract for \nxrShareNoAbstract{} of the included
papers. Work that uses none of our XR phrases is missed. The pilot contained only \nxrPilotRelevant{}
relevant papers, so precision and recall carry wide intervals. A single-label taxonomy undercounts papers
that span categories, so category sizes count the papers for which a category is \emph{primary}. Selecting
papers for deep reading by citation favours older work, and a recent-paper quota corrects this only in
part. Text-search counts depend on how authors word their abstracts, so they are indicative to within a
few points. Screening and labelling used a language model, and drafts were checked against the source
texts. Every count in this review is generated by script from the atlas data files.
